\begin{lemma}
\label{editorial-source-lemma-characterize-separable-field-extensions}
Let $k$ be a field of characteristic $p > 0$.
Let $K/k$ be a field extension.
The following are equivalent:
\begin{enumerate}
\item $K$ is separable over $k$,
\item for every $k$-linearly independent subset $\{a_1, \ldots, a_m\}$
of $K$ the set $\{a^p_1, \ldots, a_m^p\}$ is $k$-linearly independent,
\item the ring $K \otimes_k k^{1/p}$ is reduced, and
\item $K$ is geometrically reduced over $k$.
\end{enumerate}
\end{lemma}

\begin{proof}
The implication (1) $\Rightarrow$ (4) is Lemma
\ref{editorial-source-lemma-separable-extension-preserves-reducedness}
applied to every field extension of $k$ as a reduced
$k$-algebra. The symmetry map $a\otimes b\mapsto b\otimes a$
has itself as inverse. The implication (4) $\Rightarrow$ (3)
then follows by using the specified field $k^{1/p}$.

For (3) $\Rightarrow$ (2), retain the original tensor
ring $B=K\otimes_k k^{1/p}$ and the original map
$$
m:B\longrightarrow K,\qquad
\sum_i\lambda_i\otimes\mu_i\longmapsto
\sum_i\lambda_i^p\mu_i^p.
$$
This is a ring homomorphism. Additivity and multiplicativity
follow from Frobenius in characteristic $p$, and balancing
follows from
$(c\lambda)^p\mu^p=\lambda^p(c\mu)^p$ for $c\in k$.
Every $\mu_i^p$ belongs to the original $k$, so the
displayed sum belongs to $K$. For every original
$z=\sum_i\lambda_i\otimes\mu_i$ one has
$$
z^p=\sum_i\lambda_i^p\otimes\mu_i^p
=\left(\sum_i\lambda_i^p\mu_i^p\right)\otimes1
=m(z)\otimes1.
$$
Here all mixed multinomial terms have coefficients
divisible by the prime $p$. The map
$K\to B$, $a\mapsto a\otimes1$, is injective, as
extension of scalars from one field to another is
faithful on vector spaces. If $B$ is reduced, $m(z)=0$
therefore implies $z^p=0$ and then $z=0$.

Let $a_1,\ldots,a_m$ be $k$-linearly independent,
and suppose $\sum_{i=1}^m c_i a_i^p=0$ with $c_i\in k$.
Let $\mu_i\in k^{1/p}$ be the unique root with
$\mu_i^p=c_i$. Then
$m(\sum_i a_i\otimes\mu_i)=0$, so the tensor is zero.
The elements $a_i\otimes1$ are independent over
$k^{1/p}$: extend the original finite family to a
$k$-basis of $K$ and extend that basis by scalars.
Consequently every $\mu_i=0$ and every $c_i=0$.
This proves (2). The argument uses the displayed
ring homomorphism, not an assertion that $m$ is
$k$-linear: in fact $m(cz)=c^p m(z)$ for the
original scalar action of $c\in k$.

For (2) $\Rightarrow$ (1), property (2) is inherited
by every intermediate field of the original $K/k$.
It suffices, by the definition of separability, to
prove that each finitely generated such field is
separably generated. Work with one of these fields,
denoted by $K$ in the remainder of this argument.
Choose a transcendence basis $x_1,\ldots,x_d$ and
retain $K'=k(x_1,\ldots,x_d)\subseteq K$. Since
$K/k$ is finitely generated and $K/K'$ is algebraic,
it is finite by Fields, Lemma
\ref{fields-lemma-algebraic-finitely-generated}.
Among these original bases choose one for which
the positive integer $[K:K']_i$ is least. Such a
least value exists because the set of bases is
nonempty and these values are positive integers.

If $K/K'$ is separable, this basis is separating.
Otherwise choose $x_{d+1}\in K$ not separable over
$K'$, and retain the intermediate field
$E=K'(x_{d+1})$. Then $[E:K']_i>1$.
The original $E/k$ is generated by the displayed
$d+1$ elements, has $x_1,\ldots,x_d$ as a
transcendence basis, and inherits (2) from $K/k$.
Apply Lemma \ref{editorial-source-lemma-mini-separability} to $E/k$,
with its integer $n$ equal to this $d$. It supplies
$1\leq j\leq d+1$ such that
$$
K''=k(x_1,\ldots,\widehat x_j,\ldots,x_{d+1})
\subseteq E
$$
has the indicated generators as a transcendence
basis and $E/K''$ is finite separable. They are
also a transcendence basis of $K/k$, since
$K/E$ is finite algebraic. The two finite towers
are $K/E/K'$ and $K/E/K''$. Fields, Lemma
\ref{fields-lemma-multiplicativity-all-degrees}
gives the full comparison
$$
[K:K']_i=[K:E]_i[E:K']_i,
\qquad
[K:K'']_i=[K:E]_i[E:K'']_i=[K:E]_i.
$$
Thus $[K:K'']_i<[K:K']_i$, contradicting the choice
of basis. This proves separable generation of
every finitely generated intermediate field, and
hence (1) for the original extension, with no
finite-generation assumption on that extension.
When $d=0$, the same argument uses an empty basis
and $K''=k$; no omitted index is out of range.

The same original tensor maps give a stronger
description even when the equivalent conditions
fail. Choose an algebraic closure $\Omega$ of $K$
and embed $k^{1/p}$ there by its unique $p$th roots.
Put $C=Kk^{1/p}\subseteq\Omega$ and
$D=K^p k\subseteq K$, retaining both composita
with their indicated embeddings. Multiplication
defines a ring map
$$
\nu:B\longrightarrow C,\qquad
\sum_i\lambda_i\otimes\mu_i\longmapsto
\sum_i\lambda_i\mu_i.
$$
Its image is the subring of finite such sums.
For a nonzero element $w$ of this subring,
$w^p=\sum_i\lambda_i^p\mu_i^p\in K\setminus\{0\}$,
so $w^{-1}=w^{p-1}/w^p$ lies in the same subring.
It is consequently a field containing $K$ and
$k^{1/p}$, and is exactly $C$. Hence $\nu$ is
surjective. Frobenius on $C$ maps into $K$ and
sends these original sums to $\sum_i\lambda_i^p\mu_i^p$.
Its image is exactly $D$: it is a field containing
$K^p$ and $k$, and each image lies in the field
generated by them. Frobenius on a field is injective,
so it is a ring isomorphism $C\longrightarrow D$.
The original $m$ is precisely its composite with
$\nu$, followed by the inclusion $D\subseteq K$.

Since the target of $\nu$ is a field, every
nilpotent of $B$ belongs to $\ker\nu$. Conversely,
if $\nu(z)=0$, then $m(z)=\nu(z)^p=0$, and the
displayed tensor identity gives $z^p=0$. Therefore
$$
\operatorname{Nil}(B)=\ker\nu=\ker m
=\{z\in B:z^p=0\}.
$$
The exact quotient maps identify
$B/\operatorname{Nil}(B)$ with $C$ by $[z]\mapsto\nu(z)$,
and with $D$ by $[z]\mapsto m(z)$. The first is a
$K$-algebra isomorphism for the original action;
the second sends the scalar $a\in K$ to $a^p\in D$.
Neither map replaces the original tensor ring or
changes its scalar structure silently. In particular
$B$ has a unique prime ideal, namely its nilradical:
primes contain all nilpotents, and its reduced
quotient is the nonzero field $C$. The equivalent
conditions of the lemma can also be expressed as
injectivity of $\nu$, or equivalently as the original
map $K\otimes_k k^{1/p}\to Kk^{1/p}$ being an
isomorphism of $K$-algebras.

Elementwise annihilation by the $p$th power in this
description does not assert that the $p$th ideal
power of the nilradical is zero. For an exact
example, take algebraically independent $u,v$
over $\mathbf F_p$, and retain
$$
s=u^p,\quad t=v^p,\quad
k=\mathbf F_p(s,t),\quad K=\mathbf F_p(u,v).
$$
Then $k^{1/p}=K$ inside an algebraic closure of $K$.
Indeed every rational function in $s,t$ has the
unique root obtained by using the same numerator
and denominator in $u,v$, since Frobenius fixes
each coefficient in $\mathbf F_p$. The elements
$u^a v^b$ for $0\leq a,b<p$ form a $k$-basis of
$K$. Independence follows on clearing the original
denominators in any proposed relation: the exponent
pairs in its polynomial terms have distinct residues
modulo $p$. They span because powers can be reduced
using $u^p=s$, $v^p=t$, and every nonzero polynomial
denominator $q(u,v)$ has inverse
$q(u,v)^{p-1}/q(u,v)^p$, with $q(u,v)^p\in k$.

The original tensor comparison is consequently
$$
B=K\otimes_k K\simeq
K[U,V]/(U^p-s,V^p-t),\qquad
1\otimes u\longleftrightarrow U,\quad
1\otimes v\longleftrightarrow V.
$$
It fixes $K$ through the left tensor factor.
Monic division gives the basis $U^aV^b$ for
$0\leq a,b<p$ on the polynomial side; the stated
basis of the right tensor factor proves the map
and its inverse on every basis vector. Put
$\epsilon=U-u$, $\eta=V-v$. The full binomial
expansions give $\epsilon^p=U^p-u^p=0$ and
$\eta^p=V^p-v^p=0$, because every intermediate
binomial coefficient is divisible by $p$.
Evaluation at $U=u,V=v$ identifies the quotient
by $(\epsilon,\eta)$ with $K$, and all elements
of this ideal have $p$th power zero. Thus it is
the nilradical $N$. Its power $N^{2p-1}$ is zero:
each generating monomial has either epsilon or
eta exponent at least $p$. But
$\epsilon^{p-1}\eta^{p-1}\neq0$, since in the
original basis its $U^{p-1}V^{p-1}$ coefficient
is $1$ and all other terms have smaller exponent
in at least one coordinate. Hence the ideal
nilpotence index is exactly $2p-1$, and in
particular $N^p\neq0$ for every prime $p$.
\end{proof}